\chapter{符号说明与数据预处理总框架}

\section{符号说明}

\begin{table}[H]
\centering
\small
\caption{主要符号说明}
\label{tab:symbols}
\begin{tabularx}{0.92\textwidth}{C{2.2cm}X}
\toprule
符号 & 含义 \\
\midrule
$x_i,y_i$ & 问题一中校正前、校正后的位移数据，单位为mm \\
$s_t$ & $t$时刻表面位移，单位为mm \\
$v_t$ & 由位移差分得到的速度，本文主要使用6 h滚动中位速度，单位为mm/h \\
$T_1,T_2$ & 三阶段模型中的两个阶段转换节点 \\
$\tau$ & 阶段内归一化时间，$\tau\in[0,1]$ \\
$g_k(\tau)$ & 第$k$阶段的单调位移基线函数 \\
$X_t$ & 降雨、孔压、微震、爆破距离和药量等扰动特征向量 \\
$r_t$ & 多源扰动对阶段基线的残差修正量 \\
$\hat t_f$ & 由逆速度法外推得到的潜在失稳时间 \\
\bottomrule
\end{tabularx}
\end{table}

\section{数据概况}

各附件的数据规模和核心处理策略见表\ref{tab:data_overview}。附件4和附件5中爆破距离、药量存在大量空值，这类空值表示“无爆破事件”，并非传感器缺失。附件4实验集中表面位移空值为预测目标，也不能按普通缺失补齐。

\begin{table}[htbp]
\centering
\small
\caption{附件数据概况与处理策略}
\label{tab:data_overview}
\begin{tabularx}{0.98\textwidth}{C{1.4cm}C{2.8cm}C{2.3cm}X}
\toprule
附件 & 数据规模 & 主要变量 & 处理重点 \\
\midrule
1 & 10000行 & 两组位移 & 稳健仿射校准，比较不同线性稳健模型 \\
2 & 10000行 & 表面位移 & 速度去尖峰、滚动速度、复合证据阶段识别 \\
3 & 训练10000行、实验5000行 & 降雨、孔压、微震、深部/表面位移 & 缺失补齐、异常识别、变量贡献分析 \\
4 & 训练10000行、实验5000行 & 表面位移与扰动变量、阶段标签 & 阶段归一化基线、扰动残差预测、指定时刻输出 \\
5 & 9450行 & 6个候选预警变量与表面位移 & 五变量组合选择、三级预警阈值和逆速度验证 \\
\bottomrule
\end{tabularx}
\end{table}

\section{统一预处理原则}

为保证五问之间的结果可衔接，本文使用统一预处理规则。

\begin{enumerate}
  \item 时间索引统一。对含时间字段的数据直接转化为时间戳；对仅有编号的等间隔序列，按10 min间隔构造相对时间，用于速度和窗口统计。
  \item 非爆破空值置零。爆破点距离和单段最大药量的空值解释为无爆破事件，同时构造爆破冲击衰减项，避免把“无事件”当作缺失。
  \item 稀疏变量与连续变量分开处理。降雨量、微震事件数保留零值和突发峰值；孔压、深部位移、表面位移用局部趋势补齐和残差阈值识别异常。
  \item 所有模型评估采用时间块划分，避免随机打散时间序列导致未来信息泄漏。
  \item 输出结果由脚本统一生成，论文表格只摘取核心结果；完整CSV和图件保存在\texttt{code/outputs/}目录下。
\end{enumerate}

\section{特征构造框架}

对问题三至问题五，本文并不直接使用原始变量回归位移，而是构造符合监测机理的时序特征。设$r_t$为降雨量、$p_t$为孔隙水压力、$m_t$为微震事件数，则常用特征包括：
\begin{equation}
  \bar r_t^{(w)}=\frac{1}{w}\sum_{j=0}^{w-1}r_{t-j},\quad
  R_t^{(w)}=\sum_{j=0}^{w-1}r_{t-j},\quad
  \Delta p_t=p_t-p_{t-1}.
\end{equation}
其中$\bar r_t^{(w)}$表示短期平均降雨，$R_t^{(w)}$表示累积入渗效应，$\Delta p_t$刻画孔压变化率。对爆破扰动，设爆破药量为$q_t$、爆破点距离为$d_t$，构造衰减冲击项
\begin{equation}
  I_t = q_t \exp(-d_t/d_0),
\end{equation}
其中$d_0$由训练集尺度确定。非爆破时刻$q_t=0$，因此$I_t=0$。这种处理保留了“距离越远影响越弱、药量越大影响越强”的工程含义。

为避免未来信息泄漏，所有滚动特征均只使用当前及历史观测。模型验证采用连续时间块划分，而非随机划分。

\section{数据质量检查}

数据质量检查包括缺失比例、非负约束、时间间隔一致性和极端值范围四类。附件3、附件4和附件5存在较多空值，但其含义不同：附件3中的空值主要是传感器缺失，附件4实验集表面位移空值是预测目标，附件4和附件5的爆破字段空值则表示无爆破事件。只有区分这些机制，后续异常识别和预测才不会混淆“缺失”和“真实零事件”。

\section{基本假设}

\begin{enumerate}
  \item 同一附件内监测点位和工况在短期内保持一致，传感器变化主要体现为量测噪声和尺度差异。
  \item 表面位移以相对位移计，短时微小回落多由量测噪声或局部扰动导致，不改变累计变形总体增长趋势。
  \item 训练集与实验集在对应阶段内具有可比的变形规律，题目给出的实验集阶段标签可信。
  \item 预警阈值用于风险分级，不等同于确定性失稳时间；逆速度法用于辅助判断提前量和快速阶段风险。
\end{enumerate}
